%======================================================================
\section{Lamps force entropy: the sector theorem}\label{sec:sector}
%======================================================================

This section proves the lower bound of Theorem~B, and the group-theoretic half of the lower bound of Theorem~A. Its main result,
\cref{thm:sector}, is stated for an arbitrary state: if an assignment of unitaries to $p,q,t,d,x,z$ satisfies the certificate $\RP$ to
precision $e$ in a state $\rho$, nearly commutes with $\rho^{1/2}$, and keeps $J$ visibly away from $I$ in that state, then $\rho$ has
entropy at least $e^{-2+o(1)}$. \Cref{sec:game} feeds into it the models that \cref{thm:transfer} extracts from strategies, so the
theorem bounds the entanglement entropy of every near-optimal strategy. For the flat state it is the lower bound for the hyperlinear
profile of $J$ (\cref{cor:hlp-lower}). Throughout, areas are taken with respect to $\RP$ (\cref{def:RP}), $J$ also denotes the word
$xzxz$, and for $0<\eps<\frac12$ we put
\begin{equation}\label{eq:sector-bfun}
\bfun(\eps)=2^{-4(24+\log\log(1/\eps))^2}.
\end{equation}
Since $\log(1/\bfun(\eps))=O((\log\log(1/\eps))^2)$, we have $\bfun(\eps)=\eps^{o(1)}$ as $\eps\to0$, and $\bfun$ is increasing on
$(0,\frac12)$.

\paragraph{The picture.} On the first $N$ points of ray~$1$ there are $2^N$ configurations $u\subseteq\{1,\dots,N\}$, and the transporter
$T_u$ of \cref{def:transporter} carries the Pauli pair $x,z$ at the configuration $0$ to the pair $T_uxT_u^{-1},T_uzT_u^{-1}$ at $u$. In
$\GP$ these $2^N$ pairs satisfy the Pauli relations: pairs at different configurations commute, and each pair squares to the same central
involution, $(T_uxT_u^{-1}\,T_uzT_u^{-1})^2=T_uJT_u^{-1}=J$. On the sector where $J$ acts as $-1$, an exact representation therefore
carries $2^N$ anticommuting pairs, that is, $2^N$ qubits. In an approximate model the transported lamps become $2^N$ approximate Pauli
pairs measured in the sector state, and \cref{cor:count} counts them: they still force about $2^N$ bits of entropy, as long as $2^N$ times
the square of their errors stays below a constant. The errors are priced by the pair areas of \cref{thm:pair-area}. A commutator of two
transported lamps has area at most $\Lambda_N$, and in a fixed state a filling of area $\Lambda$ costs its area times the length of its
conjugators, which is again of order $\Lambda$ (\cref{lem:conjugator-length,lem:calculus}); so each error is of order
$\Lambda_N^2e/\Delta$, where $\Delta$ measures how far $J$ is from $I$. The proof balances $2^N$ against $\Lambda_N^4$; with
$\Lambda_N\le2^{(\log N+23)^2}$ the best $N$ has $2^N\approx\Delta^2e^{-2}\bfun(e)$.

\begin{theorem}[The sector theorem]\label{thm:sector}
Let $\phi:F(p,q,t,d,x,z)\to\U(D)$ be a homomorphism and $\rho$ a density matrix on $\C^D$, with $\xi=\rho^{1/2}$. Suppose that, for some
$e>0$ and $0<\Delta\le2$,
\begin{enumerate}[label=(\roman*),itemsep=1pt]
\item $\hs{(\phi(r)-I)\xi}\le e$ for every $r\in\RP$;
\item $\hs{[\phi(s),\xi]}\le e$ for every letter $s\in\{p,q,t,d,x,z\}$;
\item $\hs{(\phi(J)-I)\xi}\ge\Delta$.
\end{enumerate}
If $e\le2^{-1071}\Delta$, then
\[
\log\rank\rho\ \ge\ c_1\,\Delta^2e^{-2}\,\bfun(e),\qquad c_1=\frac1{4.1\cdot10^7}.
\]
Moreover, let $J'=\mathrm r(\phi(J))$ be the rounding of \cref{lem:rounding}, $P_\pm=\frac12(I\pm J')$, $p_-=\Tr(P_-\rho)$, and let
$\rho_-=P_-\rho P_-/p_-$ be the sector state of \cref{lem:sector-state}. Then $p_-\ge\Delta^2/16$, $\hs{[J',\xi]}\le6e$,
$S(\rho_-)\ge c_1\Delta^2e^{-2}\bfun(e)$, and
\[
S(\rho)\ \ge\ \frac{c_1}{16}\,\Delta^4e^{-2}\,\bfun(e)-1 .
\]
\end{theorem}

For the flat state $\rho=I/D$ the state commutators vanish, condition~(ii) holds trivially, and the theorem bounds $\log D$; this is
\cref{cor:hlp-lower}. Nothing about the state is used beyond (i)--(iii), so the theorem applies to any state in which a model is
accurate.

The proof follows the picture. For a scale $N$ we round the transported lamps to involutions (Step~1), derive their pairwise relations
from the pair areas and their relation to $J$ from \cref{lem:J-transport} (Step~2), pass to the sector state where the rounded $J$ is
$-1$ (Step~3), and apply the count (Step~4); Step~5 chooses $N$. Every estimate is measured in the state, and conjugators are paid for
through state commutators.

\begin{proof}
Fix $N\ge1$, to be chosen in Step~5, and put $n=\max(3,N)$, $M=2^N$, $l=21Nn\ge63$ and $\Lambda=\Lambda_N$. For $u\subseteq\{1,\dots,N\}$
the transporter $T_u$ has $\lvert T_u\rvert\le l$ (\cref{thm:pair-area}). By \cref{rem:pair-growth}, $\Lambda\ge14000N^2n=\frac{2000N}3\,l$,
so
\begin{equation}\label{eq:sector-l}
l\le\Lambda/666,\qquad l^2\le\Lambda^2/443556 .
\end{equation}
Every word of $\RP$ has length at most $30$ (in fact at most $22$, \cref{def:RP}), so \cref{lem:conjugator-length} applies with $K=30$.

\emph{Step 1: rounding.} Let $A_u=\phi(T_uxT_u^{-1})$ and $B_u=\phi(T_uzT_u^{-1})$. These words have length at most $2l+1$, so their state
commutators are at most $(2l+1)e$ by \cref{lem:calculus}(a). Since $A_u^2=\phi(T_u)\phi(x^2)\phi(T_u)^{-1}$ with $x^2\in\RP$ and
$\hs{[\phi(T_u),\xi]}\le le$, \cref{lem:calculus}(c) gives $\hs{(A_u^2-I)\xi}\le e+2le\le(2l+1)e$, and likewise for $B_u$ with
$z^2\in\RP$. Put
\[
X_u=\mathrm r(A_u),\qquad Z_u=\mathrm r(B_u),\qquad J'=\mathrm r(\phi(J)) .
\]
By \cref{lem:rounding}, $\hs{(X_u-A_u)\xi}\le(2l+1)e$ and $\hs{[X_u,\xi]}\le3(2l+1)e\le7le$, and the same holds for $Z_u$ and $B_u$. We
write
\[
\gamma=7le,\qquad\text{so that}\qquad\hs{[X_u,\xi]},\ \hs{[Z_u,\xi]}\le\gamma,\qquad \gamma\ge441e .
\]
Since $\phi(J)^2=\phi((xz)^4)$ and $(xz)^4\in\RP$, we have $\hs{(\phi(J)^2-I)\xi}\le e$, so \cref{lem:rounding} gives
$\hs{(J'-\phi(J))\xi}\le e$ and, as $J=xzxz$ has length $4$, $\hs{[J',\xi]}\le4e+2e=6e$.

\emph{Step 2: relations.} \emph{Different sites.} Let $u\ne w$, $a,b\in\{x,z\}$, and $c=[T_uaT_u^{-1},T_wbT_w^{-1}]$. This word has length
at most $4(2l+1)\le9l$ and area at most $\Lambda$ (\cref{thm:pair-area}). By \cref{lem:conjugator-length} it is freely a product of at most
$\Lambda$ conjugates of words of $\RP^{\pm1}$ with conjugators of length at most $9l+30\Lambda$, so by \cref{lem:calculus}(d)
\[
\hs{(\phi(c)-I)\xi}\ \le\ \Lambda\bigl(1+2(9l+30\Lambda)\bigr)e=\Lambda(60\Lambda+18l+1)e\ \le\ 61\Lambda^2e,
\]
because $18l+1\le\Lambda$ by~\eqref{eq:sector-l}. Let $R^a_u$ be $X_u$ if $a=x$ and $Z_u$ if $a=z$. We replace the four factors of
$\phi(c)=Y_1Y_2Y_1^{-1}Y_2^{-1}$, $Y_1=\phi(T_uaT_u^{-1})$, $Y_2=\phi(T_wbT_w^{-1})$, by the involutions $R^a_u,R^b_w,R^a_u,R^b_w$. Each
factor moves by at most $3\gamma$ in the state-weighted norm: by at most $(2l+1)e\le\gamma$ for $Y_1,Y_2$, and for the inverted factors,
writing $Y^{-1}-R=Y^{-1}(R-Y)R$, by at most $\hs{(R-Y)\xi}+2\hs{[R,\xi]}\le3\gamma$. The replacements have state commutators at most
$\gamma$, so \cref{lem:calculus}(e) bounds the total cost by $4\cdot3\gamma+2(3+2+1)\gamma=24\gamma$. Since
$[R^a_u,R^b_w]=(R^a_uR^b_wR^a_uR^b_w-I)R^b_wR^a_u$ and moving $R^b_wR^a_u$ to the right of $\xi$ costs at most $2\cdot2\gamma$,
\begin{equation}\label{eq:sector-comm}
\hs{[Z_u,Z_w]\xi},\ \ \hs{[X_u,Z_w]\xi}\ \le\ 61\Lambda^2e+28\gamma\qquad(u\ne w).
\end{equation}

\emph{One site.} Let $K_u=X_uZ_uX_uZ_u$. By \cref{lem:J-transport}, $A_uB_uA_uB_u=\phi(T_uJT_u^{-1})=\phi(j_u)\phi(J)$ with
$j_u=T_uJT_u^{-1}J^{-1}$, a word of length at most $2l+8$ and area at most $4\lvert T_u\rvert\le4l$. By
\cref{lem:conjugator-length,lem:calculus}(d),
\[
\hs{(\phi(j_u)-I)\xi}\le4l\bigl(1+2(2l+8+30\cdot4l)\bigr)e=4l(244l+17)e\le984\,l^2e .
\]
Replacing $X_u,Z_u$ by $A_u,B_u$ in $K_u$ costs at most $4(2l+1)e+2(3+2+1)(2l+1)e=16(2l+1)e\le34le$ by \cref{lem:calculus}(e), with the
state commutators of $A_u,B_u$. Moving $\phi(J)$ to the right of $\xi$ costs at most $2\cdot4e$, and replacing $\phi(J)$ by $J'$ costs at
most $e$. Hence, for every $u$,
\[
\hs{(K_u-J')\xi}\ \le\ 34le+984l^2e+8e+e\ \le\ \gamma_0:=1000\,l^2e .
\]
As $X_u$ and $Z_u$ are involutions, $X_uK_uX_u=Z_uK_uZ_u=K_u^*$, and $K_u^*-J'=K_u^*(J'-K_u)J'$, so
$\hs{(K_u^*-J')\xi}\le\gamma_0+2\hs{[J',\xi]}\le\gamma_0+12e$. Writing $X_uJ'X_u-J'=X_u(J'-K_u)X_u+(K_u^*-J')$, we get
$\hs{(X_uJ'X_u-J')\xi}\le(\gamma_0+2\gamma)+(\gamma_0+12e)$, and since $[X_u,J']=(X_uJ'X_u-J')X_u$, moving $X_u$ to the right of $\xi$ adds
$2\gamma$. With the same argument for $Z_u$,
\begin{equation}\label{eq:sector-J}
\hs{[X_u,J']\xi},\ \ \hs{[Z_u,J']\xi}\ \le\ 2\gamma_0+4\gamma+12e .
\end{equation}
Finally $\{X_u,Z_u\}=(K_u+I)Z_uX_u$. Since $(J'+I)P_-=0$,
\[
\{X_u,Z_u\}P_-=(K_u-J')Z_uX_uP_-+(J'+I)[Z_uX_u,P_-] .
\]
Applied to $\xi$, the first term has norm at most $\gamma_0+2\cdot2\gamma+2\hs{[P_-,\xi]}\le\gamma_0+4\gamma+6e$, as
$\hs{[P_-,\xi]}=\frac12\hs{[J',\xi]}\le3e$. The second has norm at most $2\hs{[Z_uX_u,P_-]\xi}=\hs{[Z_uX_u,J']\xi}$, and
$[Z_uX_u,J']=Z_u[X_u,J']+[Z_u,J']X_u$ gives the bound $2(2\gamma_0+4\gamma+12e)+2\gamma$ by~\eqref{eq:sector-J}. As $30e\le\gamma$,
\begin{equation}\label{eq:sector-anti}
\hs{\{X_u,Z_u\}P_-\xi}\ \le\ 5\gamma_0+14\gamma+30e\ \le\ 5\gamma_0+15\gamma .
\end{equation}

\emph{Step 3: the sector.} We have $p_-=\Tr(P_-\rho)=\hs{P_-\xi}^2=\frac14\hs{(I-J')\xi}^2$, and
$\hs{(I-J')\xi}\ge\hs{(I-\phi(J))\xi}-\hs{(\phi(J)-J')\xi}\ge\Delta-e\ge\Delta/2$, so $p_-\ge\Delta^2/16$. Let $\rho_-$ and
$\xi_-=\rho_-^{1/2}$ be as in \cref{lem:sector-state}; by its part~(a), $\hs{P_+\xi P_-}\le\frac12\hs{[J',\xi]}\le3e$. We bound the five
quantities that \cref{cor:count} needs, in the state $\rho_-$.
\begin{itemize}[itemsep=1pt]
\item For an operator $R$ with $\op R\le2$, \cref{lem:sector-state}(b) gives
\[
\hs{R\xi_-}\le(\hs{R\xi}+\op R\hs{[J',\xi]})/\sqrt{p_-}\le(\hs{R\xi}+12e)/\sqrt{p_-}.
\]
With~\eqref{eq:sector-comm}, $\hs{[Z_u,Z_w]\xi_-}$ and $\hs{[X_u,Z_w]\xi_-}$ are at most $(61\Lambda^2e+28\gamma+12e)/\sqrt{p_-}$.
\item By the first bound of \cref{lem:sector-state}(b) and~\eqref{eq:sector-anti},
\[
\hs{\{X_u,Z_u\}\xi_-}\le\frac{\hs{\{X_u,Z_u\}P_-\xi}+2\hs{P_+\xi P_-}}{\sqrt{p_-}}\le\frac{5\gamma_0+15\gamma+6e}{\sqrt{p_-}}.
\]
\item By \cref{lem:sector-state}(c) for the unitaries $X_u$, $Z_u$, with~\eqref{eq:sector-J},
\begin{align*}
\hs{[X_u,\xi_-]}&\le\frac{\hs{[X_u,J']\xi}+3\hs{[X_u,\xi]}+\hs{[J',\xi]}}{\sqrt{p_-}}\\
&\le\frac{(2\gamma_0+4\gamma+12e)+3\gamma+6e}{\sqrt{p_-}}=\frac{2\gamma_0+7\gamma+18e}{\sqrt{p_-}},
\end{align*}
and the same bound holds for $\hs{[Z_u,\xi_-]}$.
\end{itemize}
By~\eqref{eq:sector-l}, $\gamma=7le\le0.011\,\Lambda e$ and $\gamma_0=1000l^2e\le0.0023\,\Lambda^2e$, and $\Lambda\ge42000$. So
$61\Lambda^2e+28\gamma+12e\le(61\Lambda^2+0.31\Lambda+12)e$, $5\gamma_0+15\gamma+6e\le(0.012\Lambda^2+0.17\Lambda+6)e$ and
$2\gamma_0+7\gamma+18e\le(0.005\Lambda^2+0.08\Lambda+18)e$, and all five quantities are at most
\[
\eta:=\frac{63\Lambda^2e}{\sqrt{p_-}}\ \le\ \frac{252\Lambda^2e}{\Delta}.
\]

\emph{Step 4: counting.} Apply \cref{cor:count} on $\C^D$ to the state $\rho_-$ and the $M\ge2$ pairs of involutions $(X_u,Z_u)$,
$u\subseteq\{1,\dots,N\}$, ordered in any way; all its hypotheses hold with the uniform bound $\eta$. If
\begin{equation}\label{eq:sector-N}
2^N\Lambda_N^4\ \le\ \frac{\Delta^2}{1.02\cdot10^7\,e^2},
\end{equation}
then $16M\eta^2\le16\cdot252^2\cdot2^N\Lambda^4e^2/\Delta^2\le16\cdot252^2/(1.02\cdot10^7)<0.0997<\frac1{10}$, so \cref{cor:count} gives
$S(\rho_-)\ge\frac34M\ge2^{N-1}$. By \cref{lem:sector-entropy}, $\log\rank\rho\ge\log\rank\rho_-\ge S(\rho_-)$.

\emph{Step 5: the choice of $N$.} By \cref{rem:pair-growth}, $\log\Lambda_N\le(\log N+23)^2$. Put
\[
Q=\log\frac{\Delta^2}{1.02\cdot10^7\,e^2},\qquad g(N)=N+4(\log N+23)^2 ,
\]
so that $g$ is increasing, $g(1)=2117$, and~\eqref{eq:sector-N} holds whenever $g(N)\le Q$. For $e\le2^{-1071}\Delta$ we have
$Q\ge2142-\log(1.02\cdot10^7)>2118$, since $\log(1.02\cdot10^7)<23.3$, and $e\le\Delta/2$. Let $N$ be the largest integer with $g(N)\le Q$;
then $N\ge1$. By maximality $g(N+1)>Q$, that is, $N>Q-1-4(\log(N+1)+23)^2$. Also $N+1<g(N)\le Q<2\log(1/e)$, because $\Delta\le2$ gives
$Q\le2+2\log(1/e)-23$; so $\log(N+1)<1+\log\log(1/e)$. Therefore
\[
2^{N-1}\ >\ 2^{\,Q-2-4(\log(N+1)+23)^2}\ \ge\ \frac{\Delta^2}{4.08\cdot10^7\,e^2}\,2^{-4(24+\log\log(1/e))^2}\ \ge\ c_1\Delta^2e^{-2}\bfun(e),
\]
and with Step~4 this proves the bounds for $\log\rank\rho$ and $S(\rho_-)$. We showed $p_-\ge\Delta^2/16$ and $\hs{[J',\xi]}\le6e$ in Steps~3
and~1. Finally, by \cref{lem:sector-entropy} and $h\le1$,
\[
S(\rho)\ \ge\ p_-\,S(\rho_-)-h\bigl(\min(\hs{[J',\xi]}^2/4,\tfrac12)\bigr)\ \ge\ \frac{\Delta^2}{16}\,S(\rho_-)-1. \qedhere
\]
\end{proof}

\begin{remark}[Where the exponent $4$ comes from]\label{sector:rem-exponent}
The pair area enters~\eqref{eq:sector-N} to the fourth power because it enters the errors squared: a filling of area $\Lambda$ is charged
its area times the length of its conjugators, which is again of order $\Lambda$ (\cref{lem:conjugator-length}). This is why $\bfun$ has
the factor $4$ in its exponent. In the flat state conjugation costs nothing, the words $c$ of Step~2 have defect at most $\Lambda e$
(\cref{lem:defect}), and the same argument, with $\Lambda_N^2$ in place of $\Lambda_N^4$, gives the factor $2^{-2(24+\log\log(1/e))^2}$ for
$\log D$ (with other constants). Both factors are $e^{o(1)}$, and we do not use the sharper flat one. The exponent $2$ of $\Delta/e$ is the
same in both cases.
\end{remark}

\begin{remark}[Under the polynomial Dehn bound]\label{rem:sector-poly}
Assume~\ref{H-poly}. Then the pair areas satisfy $\Lambda_N\le C'N^\kappa$ for all $N\ge1$, with $\kappa=6+\log107$ and an absolute
constant $C'\ge1$ (\cref{rem:pair-growth}). In Step~5 choose instead $N$ maximal with $2^NC'^4N^{4\kappa}\le\Delta^2/(1.02\cdot10^7e^2)$;
such $N\ge1$ exists when $e\le e_0\Delta$ with $e_0=1/(5000\,C'^2)$, and then~\eqref{eq:sector-N} holds. Maximality gives
$2^{N-1}>\Delta^2/(4.08\cdot10^7C'^4e^2(N+1)^{4\kappa})$, and $N+1\le3\log(1/e)$ because $2^N<e^{-2}$. So, for $e\le e_0\Delta$,
\[
\log\rank\rho,\ S(\rho_-)\ \ge\ c\,\Delta^2e^{-2}(\log(1/e))^{-4\kappa},\qquad c=\frac1{4.08\cdot10^7\cdot81^{\kappa}C'^4},
\]
and the factor $\bfun(e)$ becomes a power of $\log(1/e)$, here and in everything deduced from \cref{thm:sector}.
\end{remark}

\subsection{The flat profile of \texorpdfstring{$J$}{J}}

For the flat state, \cref{thm:sector} is a statement about $\eps$-representations. We record it for every finite set of relators that
contains the certificate, and, with weaker constants, for every finite set of relators in which the certificate has bounded area. The
latter covers finite presentations of $\GP$ and of $\GPp$ on any finite generating set and any word representing $J$ (\cref{rem:GP-fp}).

\begin{corollary}[Lower bound for the profile of $J$]\label{cor:hlp-lower}
\leavevmode
\begin{enumerate}[label=(\alph*),itemsep=1pt]
\item Let $R$ be a finite set of words in $p,q,t,d,x,z$ that represent $1$ in $\GP$, with $\RP\subseteq R$. For $0<\delta\le2$ and
$0<\eps\le2^{-1071}\delta$,
\[
\log\hlp_R(J;\delta,\eps)\ \ge\ c_1\,\delta^2\eps^{-2}\,\bfun(\eps).
\]
\item Let $S$ be a finite alphabet, $R\subset F(S)$ finite and $w\in F(S)$. Suppose there are words $\hat s\in F(S)$ for
$s\in\{p,q,t,d,x,z\}$ and an integer $k\ge1$ such that, for every $r\in\RP$, the word $\hat r$ obtained from $r$ by substituting $\hat s$
for $s$ has area at most $k$ in $R$, and so has $w^{-1}(\hat x\hat z)^2$. Then for $0<\delta\le2$ and $0<\eps\le2^{-1072}\delta/k$,
\[
\log\hlp_R(w;\delta,\eps)\ \ge\ \frac{c_1}{4k^2}\,\delta^2\eps^{-2}\,\bfun(\eps).
\]
\end{enumerate}
Under~\ref{H-poly} the factor $\bfun(\eps)$ improves to $(\log(1/\eps))^{-4\kappa}$, with other constants (\cref{rem:sector-poly}).
\end{corollary}

\begin{proof}
For the flat state $\rho=I/D$ we have $\xi=D^{-1/2}I$, so $\hs{Y\xi}=\fn{Y}$ for every $Y$, every state commutator vanishes, and
$\rank\rho=D$.

(a) Let $\phi$ be an $\eps$-representation of $\langle p,q,t,d,x,z\mid R\rangle$ in $\U(D)$ with $\fn{\phi(J)-I}\ge\delta$. Conditions
(i)--(iii) of \cref{thm:sector} hold for the flat state with $e=\eps$ and $\Delta=\delta$, so $\log D\ge c_1\delta^2\eps^{-2}\bfun(\eps)$.

(b) Let $\phi$ be an $\eps$-representation of $\langle S\mid R\rangle$ in $\U(D)$ with $\fn{\phi(w)-I}\ge\delta$, and let $\phi'$ be the
homomorphism on $F(p,q,t,d,x,z)$ with $\phi'(s)=\phi(\hat s)$. By \cref{lem:defect}, $\fn{\phi'(r)-I}=\fn{\phi(\hat r)-I}\le k\eps$ for
$r\in\RP$, and $\fn{\phi'(J)-\phi(w)}=\fn{\phi(w^{-1}(\hat x\hat z)^2)-I}\le k\eps$, so $\fn{\phi'(J)-I}\ge\delta-k\eps\ge\delta/2$. Apply
\cref{thm:sector} to $\phi'$ and the flat state with $e=k\eps$ and $\Delta=\delta/2$; its hypothesis $e\le2^{-1071}\Delta$ is our
assumption on $\eps$. Since $\bfun$ is increasing and $k\eps\ge\eps$,
$\log D\ge c_1(\delta/2)^2(k\eps)^{-2}\bfun(k\eps)\ge\frac{c_1}{4k^2}\delta^2\eps^{-2}\bfun(\eps)$.
\end{proof}

When $\langle S\mid R\rangle$ presents $\GP$ (or $\GPp$) and $w$ represents $J$, the hypothesis of~(b) holds for any words $\hat s$ that
represent $p,q,t,d,x,z$, because the words $\hat r$ and $w^{-1}(\hat x\hat z)^2$ represent $1$ (\cref{lem:RP-valid}) and so have finite area; this is the
argument of \cref{lem:change-presentation}. The matching upper bound, and the state-weighted version, are \cref{thm:profile}.
